%% REGO manuscript — IEEE Transactions on Automation Science and Engineering (T-ASE)
%% Regular Paper — DOUBLE-ANONYMOUS REVIEW DRAFT
%%
%% This file is a working first draft. See PAPER_NOTES.md and PAPER_TODO.md
%% (repository root) for author metadata, AI-disclosure notes, the figure/
%% table/reference plan, and the claim-status ledger. Those files are NOT
%% part of this manuscript and must never be merged into it before the
%% double-anonymous review period ends.
%%
%% Built on the standard IEEEtran "bare_jrnl.tex" skeleton (Michael Shell,
%% IEEEtran.cls). Class and package structure preserved; explanatory
%% boilerplate comments trimmed for readability. Requires IEEEtran.cls
%% 1.8b+, standard on Overleaf's IEEE journal templates.

\documentclass[journal]{IEEEtran}

\usepackage{amsmath}
\usepackage{amssymb}
\usepackage{graphicx}
\usepackage{algorithmic}
\usepackage{array}
\usepackage{url}
\interdisplaylinepenalty=2500

\hyphenation{para-mag-net-ic reg-o-lith}

\begin{document}

% ============================================================================
% TITLE — anonymized for double-anonymous review.
% See PAPER_NOTES.md for the 8 candidate titles considered and the
% rationale for the selected one. No author names, institutions, or
% identifying \thanks entries appear below, per the T-ASE double-anonymous
% review requirement (manuscript, supplementary files, and cover letter must
% all be non-identifying).
% ============================================================================
\title{Closed-Loop Magnetic Manipulation of Free Paramagnetic Granular\\
Material for Mandrel-Free Three-Dimensional Assembly}

\author{Manuscript~submitted~for~double-anonymous~review.
\thanks{Author names, affiliations, and acknowledgments are withheld from
this version per the double-anonymous review policy of IEEE Transactions on
Automation Science and Engineering (T-ASE). See the separately maintained
author metadata sheet for the final-version author block.}}

% No \markboth author-identifying header for the anonymous version.
\markboth{IEEE TRANSACTIONS ON AUTOMATION SCIENCE AND ENGINEERING (DOUBLE-ANONYMOUS REVIEW COPY)}%
{Manuscript submitted for review}

\maketitle

% ============================================================================
\begin{abstract}
% ~200 words max, T-ASE regular paper. No citations, no equations, no
% unexplained abbreviations. Numerical claims below are limited to what is
% directly supported by the current simulation (Section VIII); no shaping
% or full-pipeline success numbers are claimed, since that validation is
% still in progress (see Limitations).
Automated assembly of loose three-dimensional structures from free granular
material, without a mold, scaffold, or mechanical mandrel, is difficult
because the actuator cannot make physical contact with every particle it
must place. We investigate whether a magnetic field, generated by a small
number of repositionable and re-strengthenable dipole sources acting only
through the real paramagnetic force law, can transport and organize free
paramagnetic granular particles into a prescribed target geometry. Because a
static magnetic field cannot hold particles in stable equilibrium across an
extended free-standing surface (a consequence of Earnshaw's theorem), the
controller must instead realize a dynamic sequence of field configurations,
with every control action expressed strictly as a change to a real
source's position, orientation, or strength rather than as an added
per-particle force. We report a discrete-element-method simulation coupling
this magnetic force model with granular contact mechanics and gravity, a
closed-loop transport controller for moving clusters of particles to target
locations, and a quantitative diagnosis of the physical mechanisms that
determine when this class of controller does and does not converge. We
characterize a specific, reproducible failure mode --- a control-law
discontinuity at a zone-switching boundary --- and report cluster-integrity
and force-budget evidence supporting the physical plausibility of the
approach. [PLACEHOLDER: one-sentence summary of final shaping-stage result,
pending Section VIII.]
\end{abstract}

\begin{IEEEkeywords}
Automation, magnetic manipulation, granular materials, particle assembly,
magnetic actuation, closed-loop control, discrete element method, granular
robotics, in-situ resource utilization, extraterrestrial construction.
\end{IEEEkeywords}

% ============================================================================
% NOTE TO PRACTITIONERS — mandatory for T-ASE, 100-300 words, distinct from
% the Abstract, written for a practicing engineer outside this specialty.
% IEEEtran has no dedicated environment for this; T-ASE convention is a
% labeled paragraph immediately following the keywords/before Section I.
% ============================================================================
\section*{Note to Practitioners}
Assembling loose granular material into a three-dimensional shape normally
requires a mold, a robotic arm that touches each particle, or a container
that constrains the final geometry. In environments where transporting or
maintaining such tooling is expensive or impossible --- for example,
building structures from loose soil on the Moon, where every kilogram of
hardware is costly to deliver --- a method that shapes material using only
an external field, with no mechanical contact and no physical form, would be
valuable. This paper studies whether a small number of movable
electromagnet-like sources can play this role for a granular material that
responds weakly to magnetic fields (paramagnetic regolith simulant). We show
that the field sources must be actively moved and re-strengthened in a
planned sequence, rather than held in one fixed configuration, because a
fixed field cannot hold particles spread across an open surface. We report a
working closed-loop control scheme for moving clusters of particles to
target locations and a precise, physically grounded account of where and
why the present controller does not yet reliably bring every batch of
particles fully to rest at its target. Practitioners considering
magnetic or other field-based manipulation of granular or particulate
material --- in manufacturing, additive construction, or automated material
handling --- should take from this work that (1) actuator motion planning,
not merely field strength, determines whether an extended target shape is
reachable, and (2) the cross-talk between an active actuator and
already-placed material must be budgeted for explicitly, since it is a
principal source of failure in early designs. The system reported here is a
simulation study; no physical hardware has been built or tested.

\IEEEpeerreviewmaketitle

% ============================================================================
\section{Introduction}
% VERIFIED framing only. No fabricated motivation numbers.
\IEEEPARstart{A}{utomating} the assembly of a three-dimensional structure
from loose, uncontained granular material is a qualitatively different
problem from robotic assembly of discrete rigid parts. There is no single
object to grasp, no fixed set of mating features, and---if a mold or
mechanical scaffold is disallowed---no physical boundary condition to lean
on. This problem is of direct interest for automation in environments where
delivering and maintaining tooling is itself the dominant cost or risk: for
example, using an external field to organize loose regolith into a
structural form during lunar or planetary construction would avoid
transporting a mold, formwork, or a manipulator capable of contacting every
grain.

This paper investigates whether a magnetic field, produced by a small number
of external dipole-like sources whose position, orientation, and strength
are the only control inputs, can (a) transport discrete clusters of
paramagnetic granular particles from an initial distribution to prescribed
target locations, and (b) reorganize those clusters into an extended target
surface, using only the real paramagnetic force law acting on free
particles --- with no mechanical mandrel, no per-particle force filtered by
a particle's intended destination, and no force term invented to make a
target shape reachable.

A central physical constraint shapes the entire controller design. In a
current-free region, $\nabla^2(B^2) = 2|\nabla B|^2 \geq 0$, so $B^2$ is
subharmonic and its \emph{static} maxima occur only at the field sources
themselves (Earnshaw's theorem \cite{earnshaw_wiki}). Consequently, no
static field configuration can hold paramagnetic particles in stable
equilibrium while they are spread across an extended free-standing surface;
every such particle collapses toward the nearest source. Any mechanism that
appears to hold particles on a static extended surface must therefore
either be genuinely time-varying (with the source revisited at a rate fast
relative to the mechanical response of the particles
\cite{earnshaw_timeaveraging_patent}) or rely on a real secondary physical
mechanism such as granular contact and gravity. This paper takes that
constraint as a design requirement rather than an obstacle to work around:
every control action described here changes a real source's configuration,
never a per-particle force law.

\subsection{Contributions}
% Derived from what is ACTUALLY implemented/validated in this session and
% the codebase, not aspirational. Update once shaping is finalized.
This paper makes the following contributions:
\begin{enumerate}
\item A physically constrained simulation coupling an analytical
paramagnetic dipole force model with granular contact mechanics and
gravity, in which every control action is realized strictly as a change to
a real source, never as an added per-particle force (Section~IV).
\item A closed-loop transport controller that moves clusters of free
particles to target locations using real position-and-velocity feedback,
replacing an earlier open-loop, time-scheduled design (Section~V).
\item A quantitative diagnosis, verified against real simulation state
rather than inferred from visualization, of the specific control-law
mechanisms that cause a discrete-time realization of this controller to
fail to converge, including a projection/rate-limiting interaction and a
zone-switching discontinuity (Section~VI).
\item Cluster-integrity and force-budget evaluation metrics --- centroid,
RMS spread, maximum particle radius, saturation fraction --- that
distinguish ``the centroid reached the target'' from ``the material
remained physically coherent while being moved'' (Section~VII).
\item [PLACEHOLDER: shaping-stage contribution, pending validation of the
coverage-feedback traveling-well design in Section~VI-D.]
\end{enumerate}

The remainder of this paper is organized as follows. Section~II reviews
related work. Section~III formalizes the problem. Section~IV presents the
physical model. Section~V describes the transport controller. Section~VI
reports the controller-design findings and the current state of the
shaping problem. Section~VII defines the evaluation methodology.
Section~VIII reports results. Section~IX discusses the findings, and
Section~X states limitations.

% ============================================================================
\section{Related Work}
% Real citations only for sources actually found and read this session
% (Earnshaw time-averaging, moving-trap decelerator, magnetophoretic
% conveyor, dynamically-changing-landscape substrate trajectories). All
% other subsections are left as clearly marked placeholders rather than
% populated with invented references. See PAPER_TODO.md "Reference audit."

\subsection{Automated Granular and Particulate Assembly}
[PLACEHOLDER --- NEEDS LITERATURE: granular/particulate automated assembly,
granular additive manufacturing, particle-based construction. No sources
verified yet; do not cite until confirmed real.]

\subsection{Magnetic Particle Manipulation}
Moving magnetic traps have been used to perform real, controlled work on a
particle ensemble --- for example, a translating magnetic well was used to
decelerate a beam of paramagnetic atoms, removing over 98\% of its initial
kinetic energy \cite{moving_trap_decelerator} --- demonstrating that a
\emph{translating}, not static, potential well is a physically established
actuation primitive. Magnetophoretic conveyors built from micro-stripe or
electromagnet arrays, switched in a periodic sequence, create a traveling
potential well that transports and redistributes particle populations along
a track \cite{magnetophoretic_conveyor_patent}, and dynamically changing
magnetic field landscapes have been used to prescribe close-to-substrate
particle trajectories \cite{substrate_trajectories}. [PLACEHOLDER --- NEEDS
LITERATURE: magnetic tweezers, colloidal magnetic assembly, comparison of
which of these results transfer from single-particle/colloidal/fluid
systems to a dry granular ensemble.]

\subsection{Magnetic Field Control and Optimization}
[PLACEHOLDER --- NEEDS LITERATURE: feedback control of magnetic
manipulators, trajectory tracking, model-predictive magnetic control,
multi-source field optimization. Distinguish single-object trajectory
tracking (the majority of this literature) from the ensemble
coverage/redistribution objective this paper requires.]

\subsection{Extended Static Magnetic Equilibria and Earnshaw's Theorem}
Earnshaw's theorem is well established for electrostatic and magnetostatic
configurations of point sources \cite{earnshaw_wiki}; the standard
work-around is time-averaging a periodic field at a frequency large
compared to the growth rate of the resulting instability
\cite{earnshaw_timeaveraging_patent}. This paper applies that principle at
the level of a discretely re-positioned single source rather than a
continuously oscillating field, and reports (Section~VI) direct simulation
evidence, not just the theoretical argument, for why a static extended
attractor was not viable for this system.

\subsection{Lunar and Extraterrestrial Regolith Handling}
[PLACEHOLDER --- NEEDS LITERATURE: lunar regolith mechanics, in-situ
resource utilization construction concepts, automated extraterrestrial
assembly. Motivation only; this paper does not model lunar vacuum,
thermal, or radiation effects (Section~X).]

\subsection{Research Gap}
Existing magnetic manipulation and trajectory-control literature
predominantly addresses single objects, microrobots, or particles suspended
in a fluid, where drag dominates the dynamics and there is no dry contact
network. Existing granular assembly and additive-manufacturing literature
predominantly relies on physical confinement (a mold, a substrate, a
build plate) to fix the final geometry. Neither directly addresses a dry,
gravity-loaded, contact-dominated granular ensemble that must be organized
into a free-standing extended target surface using only a field-based
actuator with no mechanical boundary condition. This paper addresses that
gap.

% ============================================================================
\section{Problem Formulation}
We consider $N$ paramagnetic spherical particles of radius $R$, density
$\rho$, and volumetric magnetic susceptibility $\chi$, initially distributed
on a floor under a uniform gravitational acceleration $g$. A set of $N_d$
magnetic dipole sources, each with a controllable position
$\mathbf{p}_k(t)$, moment direction, and strength $s_k(t)\in[0,1]$, produces
a field $\mathbf{B}(\mathbf{x},t)=\sum_k \mathbf{B}_k(\mathbf{x},t)$. Every
particle experiences the same force law (Section~IV-B) regardless of which
source(s) are ``intended'' for it --- the simulation and controller never
filter forces by a particle's assigned cluster or destination.

The overall task is divided into phases with distinct objectives, which we
keep explicitly separate rather than folding into one control objective:

\begin{itemize}
\item \textbf{Clustering}: gather $N$ particles from an initial spread
distribution into $K$ coherent groups of known, permanent identity.
\item \textbf{Transport}: move each group's centroid to a prescribed target
region while keeping the group's real-time position error and velocity
below defined tolerances (Section~V), without the mechanism used to move
one group perturbing another group's already-achieved state.
\item \textbf{Shaping}: reorganize a group, once delivered, from a compact
cluster into a spatial distribution over a prescribed extended target
surface (a disk or a cylindrical band), subject to the surface-coverage
and density-uniformity objective formalized in Section~VI-D, not a
centroid-only objective.
\item \textbf{Holding}: maintain the shaped distribution for a defined
duration without a static-equilibrium assumption that Section~III's
Earnshaw argument already rules out.
\end{itemize}

% ============================================================================
\section{Physical Model}

\subsection{Particle Dynamics}
Each particle $i$ obeys
\begin{equation}
m_i\ddot{\mathbf{x}}_i = \mathbf{F}_{\mathrm{mag},i} + \mathbf{F}_{\mathrm{grav},i}
+ \mathbf{F}_{\mathrm{contact},i},
\label{eq:eom}
\end{equation}
with $m_i = \rho V_p$, $V_p = \tfrac{4}{3}\pi R^3$. Gravity is applied
unconditionally every integration step,
$\mathbf{F}_{\mathrm{grav},i} = -m_i g\,\hat{\mathbf{z}}$, independent of any
magnetic actuation; this separation is deliberate (Section~V-B) so that
every controller correction is expressed relative to a force budget that
already includes gravity, not as an ad hoc compensation term.

\subsection{Paramagnetic Force Law}
Each particle experiences
\begin{equation}
\mathbf{F}_{\mathrm{mag},i} = \frac{V_p\,\chi_{\mathrm{eff}}(B_i)}{2\mu_0}\,
\nabla\!\left(B^2\right)\Big|_{\mathbf{x}_i},
\label{eq:parforce}
\end{equation}
where $\chi_{\mathrm{eff}}$ is a field-dependent effective susceptibility
that saturates at high field magnitude, and $\nabla(B^2)$ is computed
analytically (not by finite difference) from the superposition of all
active dipole sources' fields. A soft, smoothly saturating clamp bounds
$|\nabla(B^2)|$ so that near-field forces remain finite and the value at
which the clamp engages is verified, not assumed, against the standoffs
actually used by the controller (Section~V).

\subsection{Contact Mechanics}
Particle--particle and particle--boundary contacts use a Hertzian normal
contact law with velocity-proportional damping and a Coulomb-limited
tangential friction term. The integration timestep is chosen so that
$\omega\,\Delta t$ for the stiffest contact realized in a given run remains
within the stability margin, checked per run rather than assumed from a
default parameter set (see \texttt{analysis/validate\_phase2.py} in the
accompanying code).

\subsection{Physical Constraints and Known Non-Physical Placeholders}
As established in Section~I, no static configuration of the sources in
\eqref{eq:parforce} can produce a stable equilibrium for particles
distributed across an extended free surface. One implementation-level
placeholder inherited from an earlier development stage,
\texttt{surf\_conf}, applies a linear spring toward a target surface for
particles during the shaping phase; it has no physical source (vacuum
exerts no restoring force) and is explicitly not extended, retuned, or
cited as evidence that magnetic shaping works on its own --- it stands in
for whatever a finalized dynamic control law must actually provide, and its
removal/replacement is the subject of Section~VI-D. No other artificial
per-particle force is present in the simulation.

% ============================================================================
\section{Transport Controller}

\subsection{Design Principle: Feedback Control of the Source}
Every controller decision changes only a real source's position,
orientation, and strength; the force law \eqref{eq:parforce} is never
modified, filtered by particle or cluster identity, or supplemented. This
mirrors how a real closed-loop electromagnet (a magnetic bearing, a
magnetic-tweezer instrument) is actually driven: a controller reads a
position/velocity sensor and adjusts the actuator, and the physics the
object feels is unchanged.

\subsection{Single-Dipole Vector Realization}
A single on-axis dipole can deliver a force of any commanded direction and
magnitude, up to the ceiling set by its standoff, by placing it at
$\mathbf{x}_{\mathrm{cluster}} + r\hat{\mathbf{a}}$ with moment along
$-\hat{\mathbf{a}}$, where $\hat{\mathbf{a}}$ is the desired force
direction and $r$ is a fixed design standoff. Every phase below computes a
desired acceleration vector $\mathbf{a}_{\mathrm{cmd}}$ and realizes it
through this mechanism.

\subsection{State Machine}
The controller advances a cluster through
LIFTOFF~$\rightarrow$~LIFT~$\rightarrow$~CRUISE~$\rightarrow$~BRAKE~$\rightarrow$~SETTLE.
CRUISE uses two independent deadbeat-style channels (vertical: a
signed, distance-scaled reference speed replacing an earlier discontinuous
ceiling-only law; horizontal: pure-pursuit along a precomputed,
collision-avoiding waypoint path) composed into one 3-D command and
realized per Section~V-B. BRAKE uses a critically-damped state-feedback
law,
\begin{equation}
\mathbf{a}_{\mathrm{net}} = -2\zeta\omega_n\mathbf{v}
+ \omega_n^2(\mathbf{x}_{\mathrm{target}}-\mathbf{x}),
\quad \zeta=1,
\label{eq:brake}
\end{equation}
replacing an earlier exact kinematic stopping-distance law
$a_{\mathrm{needed}}=v^2/(2d_{\mathrm{remaining}})$ that proved to have no
damping margin under discrete-time control (Section~VI-B). Arrival is
declared only once position error $d<\varepsilon_x$ and speed
$v<\varepsilon_v$ hold simultaneously for a sustained dwell, not on a fixed
timer; $\varepsilon_x=5R$ (sub-particle-radius precision is not
meaningful) and $\varepsilon_v=\varepsilon_x/\Delta t_{\mathrm{ctrl}}$ (the
speed below which a cluster cannot leave the position tolerance even
coasting, unpowered, for one control cycle). A diagnostic stall safety net
force-completes a transport that has not converged within a generous
multiple of the expected transit time, logging the event explicitly rather
than treating it as ordinary completion.

\subsection{Direction-Change Rate Limiting}
Because the dipole's standoff is fixed and only its aim direction
$\hat{\mathbf{a}}$ varies, an unconstrained per-step direction change can
sweep the dipole close enough to individual particles to produce
unphysically large transient forces. A rate limiter caps the per-step
direction change to $\Delta\theta_{\max}$, derived from the chord-to-angle
relation at the design standoff and the real, measured cluster extent
(Section~VI-A).

% ============================================================================
\section{Controller-Design Findings}
% Verified, real findings from this session's F19-F24 chain and the
% shaping audit. This section reports diagnosed mechanisms, not just fixes.

\subsection{Finding: Zone-Switching Chatter and the Realization Dead Zone}
An early implementation of the rate limiter above (Section~V-D) combined
with a magnitude-realization rule that clips the commanded force to zero
whenever the rate-limited direction deviates by more than $90^\circ$ from
the raw desired direction. When a velocity-direction reversal is
required---exactly the situation at a BRAKE-zone entry following an
overshoot---this produces a multi-control-step window of zero commanded
thrust, during which the cluster coasts unopposed, degrading whatever
braking distance estimate the controller was relying on and producing a
larger excursion than the one being corrected. This was diagnosed from
real per-control-step simulation state (not code inspection alone) and
confirmed at a second, independent near-arrival event without launching a
new simulation.

\subsection{Finding: Deadbeat Braking Has No Discrete-Time Damping Margin}
The exact kinematic stopping-distance law
$a_{\mathrm{needed}}=v^2/(2d_{\mathrm{remaining}})$, recomputed fresh every
6\,ms control step, assumes continuous re-evaluation lands velocity at
exactly zero at exactly the remaining distance --- an assumption that only
holds in continuous time. Under real discrete-time sampling, it has no
damping margin, so a re-entry into the BRAKE zone at a nontrivial incoming
speed can chatter between the BRAKE and CRUISE zones rather than converge.
Replacing this law with the critically-damped state-feedback law
\eqref{eq:brake} reduced the number of zone re-entries on the worst
observed case from 4 to 2 and reduced settling time on a synthetic
worst-case scenario from 1.21\,s to 0.33\,s, at more than $15\times$ margin
below the actuator's saturated force ceiling, but did not by itself achieve
full convergence: real closed-loop transport runs continued to stall short
of the arrival criterion, with the largest remaining velocity excursions
traced to the CRUISE$\rightarrow$BRAKE hand-off itself --- the two zones'
control laws are two different vector fields with no continuity guarantee
where they meet at the switching threshold. Transport is therefore reported
in this draft as a characterized, not fully converged, closed-loop system;
see Section~X.

\subsection{Finding: Undefended Clusters Are Physically Draggable by Other Active Sources}
% NEW, session finding. Real, verified via trajectory reconstruction.
Because the force law \eqref{eq:parforce} is deliberately not filtered by
particle or cluster identity (Section~V-A), a cluster that is not currently
being actively controlled --- either because it is still waiting for its
own transport turn, or because it was marked complete while its hold
actuator was too far from it to exert a meaningful restoring force ---
has no protection against a different, currently-active source. Real
archived trajectory data traced an event in which a transport actuator's
own large excursion (Section~VI-A/B) became, transiently, the dominant
field for two other, undefended clusters, physically drawing all
particles from multiple clusters together before the shaping phase began.
This is reported as a real, physically explicable consequence of the
no-filtering design choice interacting with an unconverged transport
controller, not a data-corruption or bookkeeping artifact (particle
identity and count remained intact throughout).

\subsection{Shaping: Root Causes and Candidate Architecture (In Progress)}
Independent of Section~VI-C, an audit of the shaping-phase implementation
identified three further, code-verified mechanisms: (i) the cap-spreading
mechanism relies on a single contact-driven compression burst with no
sustained outward driving force once particles jam; (ii) the curved-wall
mechanism provides only radial confinement toward the target radius, with
no sustained vertical or angular confinement, so particles are not held
within their intended axial band once their brief active-scan interval
ends; (iii) an existing, partially mitigated instance of source-to-source
cross-talk during the wall-scan sub-phase. Consistent with the Earnshaw
argument of Section~I, no static extended configuration is expected to
resolve any of these three. A coverage-feedback traveling-well
architecture --- one active source per cluster, continuously
present rather than only during a fixed time slot, whose position is
driven by real-time surface-occupancy feedback rather than a
pre-scripted path --- is the current candidate design, chosen over a
discrete sequential-patch-deposition alternative and a rotating
multi-source alternative on implementation-cost and cross-talk grounds
(see project development notes). \textbf{This design has not yet been
implemented or validated}; Sections~VII--VIII report only the transport
evaluation methodology and results, and shaping validation is marked
pending throughout.

% ============================================================================
\section{Evaluation Methodology}

\subsection{Transport Acceptance Criterion}
A transport attempt is judged converged only if it reaches a genuine,
criterion-based arrival (Section~V-C) rather than a stall-forced
completion, while maintaining cluster integrity throughout (below).
Reaching the target centroid position is necessary but not sufficient.

\subsection{Cluster-Integrity Metrics}
Centroid position alone cannot distinguish ``the centroid reached the
target'' from ``the material remained physically coherent while being
moved.'' For a cluster with particle positions $\mathbf{x}_i$ and centroid
$\bar{\mathbf{x}}=\tfrac{1}{N}\sum_i\mathbf{x}_i$, we report:
\begin{align}
R_{\mathrm{RMS}} &= \sqrt{\tfrac{1}{N}\textstyle\sum_i \|\mathbf{x}_i-\bar{\mathbf{x}}\|^2}, \label{eq:rms}\\
R_{\max} &= \max_i \|\mathbf{x}_i-\bar{\mathbf{x}}\|, \label{eq:rmax}\\
R_{95} &= \text{95th percentile of } \|\mathbf{x}_i-\bar{\mathbf{x}}\|, \label{eq:r95}
\end{align}
each normalized against a baseline captured immediately before that
cluster's transport attempt begins.

\subsection{Force-Budget and Saturation Metrics}
Peak commanded acceleration, peak realized force-to-weight ratio, fraction
of control steps at which the near-field clamp saturates, and minimum
source-to-particle separation are reported for every transport attempt, to
verify that the actuator was never operating at an unphysical or
unverified extreme.

\subsection{Shaping Metrics (Formulated, Pending Application)}
% Section 17 requirements, formalized now for later use.
For a target surface parameterized in cylindrical coordinates
$(r,\theta,z)$ (curved wall, $r\approx R_{\mathrm{cyl}}$) or as a target
plane (caps), we define per-particle surface-distance error
$e_{r,i}=|r_i-R_{\mathrm{cyl}}|$ (wall) or perpendicular distance to the
target plane (caps); surface coverage as the fraction of a binned
$(\theta,z)$ or $(r,\theta)$ grid occupied by at least one particle within
tolerance; density uniformity as the coefficient of variation of
per-bin occupancy; escape fraction as the fraction of particles outside a
defined tolerance of the target surface; and particle conservation as the
fraction of the original $N$ particles retained and trackable throughout.
These metrics are defined here so they can be applied without redefinition
once the shaping architecture of Section~VI-D is implemented and run.

% ============================================================================
\section{Results}
% NO FABRICATED NUMBERS. Explicit placeholders only, per project policy.
[PLACEHOLDER --- NEEDS DATA: Table~I, per-cluster transport metrics
(arrival error, final speed, transport duration, peak velocity, peak
acceleration, force/weight ratio, saturation fraction, minimum
separation, $R_{\mathrm{RMS}}$/$R_{\max}$/$R_{95}$ and their ratios to the
pre-transport baseline, pass/fail against Section~VII-A) for all four
transported clusters under the current (critically-damped BRAKE)
controller.]

[PLACEHOLDER --- NEEDS DATA: Fig.~4, representative real transport
trajectory ($z$, $d$, $v$ vs.\ $t$) for the cluster with the clearest
zone-switching event described in Section~VI-B.]

[PLACEHOLDER --- NEEDS DATA: shaping-stage coverage, density-uniformity,
and surface-distance-error results, pending Section~VI-D implementation
and the isolated shaping-fixture validation described in project
development notes.]

% ============================================================================
\section{Discussion}
The controller-design findings in Section~VI are, at present, the paper's
most load-bearing result: they identify, with evidence traced to real
per-control-step simulation state rather than inferred from an animation,
\emph{why} a physically careful, incrementally-corrected closed-loop
controller of this class fails to fully converge, and they distinguish a
genuinely fixed mechanism (the near-field realization dead zone, the
deadbeat braking instability) from a still-open one (the zone-switching
discontinuity; cross-cluster leakage under an unconverged transport
controller). [PLACEHOLDER --- expand once Section~VIII data are final:
quantitative comparison against the force/margin budget of Section~V,
and whether the remaining transport limitation and the shaping root
causes of Section~VI-D are best understood as one connected control-design
problem or two separable ones.]

% ============================================================================
\section{Limitations}
% Real, current limitations. No overclaiming.
\begin{itemize}
\item The transport controller does not yet reliably reach the strict
arrival/dwell criterion of Section~V-C for all clusters; the specific
remaining mechanism (a control-law discontinuity at the CRUISE/BRAKE
switching boundary) is characterized in Section~VI-B, not hidden.
\item The shaping-stage controller redesign of Section~VI-D is a proposed,
not yet implemented or validated, architecture.
\item The current particle count (256 total, 64 per cluster) is a
qualitative validation scale, not a field-representative granular volume;
whether the reported mechanisms hold at larger particle counts is untested.
\item The paramagnetic force model uses an on-axis dipole approximation for
control-law \emph{design} decisions (standoff, strength solving); the
physics engine's own force computation uses the full off-axis analytical
field, cross-checked against but not identical to the on-axis design
estimates in general geometry.
\item This is a simulation-only study; no physical hardware demonstration
has been performed.
\item Lunar vacuum, thermal cycling, dust adhesion/cohesion, and radiation
effects on real regolith are outside this model's scope; only bulk gravity
and Hertzian dry contact mechanics are represented.
\end{itemize}

% ============================================================================
\section{Conclusion}
[PLACEHOLDER --- to be finalized once Sections VIII-IX are complete. Draft
framing: this paper reports a physically constrained magnetic granular
manipulation simulation, a closed-loop transport controller with a
precisely diagnosed convergence limitation, and a shaping-stage redesign
motivated by a first-principles Earnshaw argument and independent
simulation evidence, with shaping-stage validation as the immediate next
step rather than a claimed result.]

% ============================================================================
\appendices
\section{Dipole Force Realization}
[PLACEHOLDER: closed-form inversion used to solve source strength for a
commanded acceleration magnitude at a fixed standoff, referenced from
Section~V-B.]

% No \section*{Acknowledgment} in the anonymous version --- funding and
% acknowledgment text can identify authors/institutions and must be
% withheld until after review, per T-ASE double-anonymous policy. See
% PAPER_NOTES.md for the final-version acknowledgment text placeholder.

% ============================================================================
% REFERENCES — only sources actually located and read this session are
% included as real \bibitem entries. Every other citation used above is a
% bracketed [PLACEHOLDER] in text, not a fabricated numbered reference.
% See PAPER_TODO.md "Reference audit" for verification status of each.
% ============================================================================
\begin{thebibliography}{9}

\bibitem{earnshaw_wiki}
``Earnshaw's theorem,'' \emph{Wikipedia, The Free Encyclopedia}. [Online].
Available: https://en.wikipedia.org/wiki/Earnshaw's\_theorem
[PROVISIONAL --- replace with a primary/archival source before submission;
an encyclopedia entry is not an acceptable final citation.]

\bibitem{earnshaw_timeaveraging_patent}
[PROVISIONAL --- passive magnetic bearing / stabilizer patent describing
time-averaging of a periodic magnetic field to circumvent Earnshaw's
theorem, located via search this session. Needs full bibliographic
verification (patent number, assignee, year) before citation.]

\bibitem{moving_trap_decelerator}
[PROVISIONAL --- ``A Moving Magnetic Trap Decelerator: A New Source for
Cold Atoms and Molecules,'' arXiv preprint, located this session. Needs
author list, arXiv ID, and eventual publication venue verified before
citation.]

\bibitem{magnetophoretic_conveyor_patent}
[PROVISIONAL --- ``Device and method for particle complex handling''
(magnetophoretic conveyor via microstripe/electromagnet array), patent
family located this session. Needs patent number and assignee verified
before citation.]

\bibitem{substrate_trajectories}
[PROVISIONAL --- ``Three-dimensional close-to-substrate trajectories of
magnetic microparticles in dynamically changing magnetic field
landscapes,'' arXiv preprint, located this session. Needs author list and
arXiv ID verified before citation.]

\end{thebibliography}

\end{document}
